\documentclass[11pt,a4paper]{article}

% ============================================
% PACKAGES
% ============================================
\usepackage[utf8]{inputenc}
\usepackage[T1]{fontenc}
\usepackage{lmodern}
\usepackage[margin=2.5cm]{geometry}
\usepackage{amsmath,amssymb,amsthm}
\usepackage{graphicx}
\usepackage{booktabs}
\usepackage{array}
\usepackage{enumitem}
\usepackage{xcolor}
\usepackage{hyperref}
\usepackage{cleveref}
\usepackage{tikz}
\usetikzlibrary{arrows.meta,positioning,shapes,calc,decorations.pathreplacing}
\usepackage{natbib}
\usepackage{caption}
\usepackage{float}
\usepackage{orcidlink}

% ============================================
% CUSTOM COLORS (consistent with Paper 1)
% ============================================
\definecolor{exitcolor}{RGB}{220,53,69}
\definecolor{codecolor}{RGB}{40,167,69}
\definecolor{auditcolor}{RGB}{0,123,255}
\definecolor{governcolor}{RGB}{255,193,7}
\definecolor{forkcolor}{RGB}{111,66,193}

% ============================================
% THEOREM ENVIRONMENTS
% ============================================
\theoremstyle{definition}
\newtheorem{definition}{Definition}[section]
\newtheorem{proposition}{Proposition}[section]
\newtheorem{principle}{Principle}[section]
\newtheorem{claim}{Claim}[section]

% ============================================
% CUSTOM COMMANDS
% ============================================
\newcommand{\papertitle}{Epistemic Capture and the Action Boundary: Corrigibility for Learned and Agentic Systems}

% ============================================
% METADATA
% ============================================
\title{\papertitle}
\author{
  Anivar A Aravind~\orcidlink{0009-0009-8995-0005} \\
  Independent Researcher\\
  Bengaluru, India \\
  \texttt{ping@anivar.net}
}
\date{February 2026 (Working Draft)}

% ============================================
% KEYWORDS ENVIRONMENT
% ============================================
\newenvironment{keywords}{\paragraph{Keywords:}}{}

% ============================================
% DOCUMENT
% ============================================
\begin{document}

\maketitle

\begin{abstract}
As digital infrastructure increasingly incorporates machine learning and agentic systems, deterministic rule execution is replaced by probabilistic inference. This shift transforms digital public infrastructure into epistemic public infrastructure (EPI), where system outputs constitute claims about reality rather than mere transactional processing. Corrigibility in such systems requires additional structural guarantees beyond those sufficient for deterministic architectures.

This paper extends a structural theory of corrigibility to learned and agentic systems. It argues that transparency in EPI must encompass logic, weights, data provenance, and representational schema (LWD-R) to enable meaningful reconstruction of decision boundaries. It further identifies compute capture as a structural constraint on practical forkability, variety drift as a temporal degradation of epistemic adequacy under distributional shift, and the action boundary protocol as a necessary separation between probabilistic inference and deterministic execution.

The analysis distinguishes structural corrigibility from epistemic competence: while corrigibility ensures that systemic error can be contested and corrected, it does not guarantee scientific or predictive optimality absent engineered constraint mechanisms. The paper proposes a layered architecture for bounding epistemic power in agentic systems and clarifies the structural conditions under which learned infrastructure remains publicly corrigible.
\end{abstract}

\begin{keywords}
Epistemic Public Infrastructure, AI Governance, Corrigibility, Action Boundary Protocol, Compute Capture, LWD-R, Workflow Capture, Agentic Skills
\end{keywords}

\paragraph{Scope.} This paper extends a structural theory of corrigibility to learned and agentic systems. It does not claim that corrigibility ensures optimal AI performance, alignment with operator intent, or epistemic accuracy. The claim is narrower: absent structural corrigibility, error in learned systems cannot be durably corrected by affected participants.

\section*{Executive Summary}

As digital systems increasingly rely on machine learning, infrastructure no longer merely processes transactions. It generates classifications and claims about reality. This shift creates \textbf{epistemic public infrastructure (EPI)}, where decision-making authority is embedded in learned models.

Extending a structural theory of corrigibility to learned systems, this paper argues that traditional transparency measures are insufficient. Meaningful reconstruction of model behavior requires disclosure of logic, weights, data provenance, and representational schema (\textbf{LWD-R}). Even when transparency is present, practical forkability may be constrained by compute asymmetry (``\textbf{compute capture}''), and performance may degrade over time under distributional shift (``\textbf{variety drift}'').

The paper introduces the \textbf{action boundary protocol}: probabilistic inference must be separated from deterministic execution through enforceable validation layers. Without such boundaries, agentic systems can exercise unbounded epistemic power.

Corrigibility ensures that learned infrastructure can be contested and corrected. It does not guarantee epistemic competence. Sustainable governance of AI-mediated infrastructure requires both structural feedback mechanisms and engineered constraint systems that function as ``compilers for reality.''

\section{Introduction}
\label{sec:intro}

The companion paper \citep{paper1} establishes a corrigibility framework for Digital Public Infrastructure (DPI), deriving five structural tests from cybernetics, commons governance, and free software. That framework addresses \textbf{deterministic} infrastructure: ledgers, registries, and payment rails where source code is the arbiter of behavior. This paper extends the framework to \textbf{stochastic} infrastructure, what we term \textbf{Epistemic Public Infrastructure (EPI)}: systems where state decisions emerge from learned parameters rather than explicit logic.

\begin{definition}[Epistemic Public Infrastructure]
\label{def:epi}
EPI comprises AI systems that determine citizen outcomes at population scale, including: identity verification models, welfare eligibility systems, public-facing chatbots with administrative authority, and autonomous agents operating within mandatory infrastructure.
\end{definition}

\subsection{DPI and EPI: Structural Contrast}
\label{subsec:dpi-epi-contrast}

Digital Public Infrastructure (DPI) processes transactions through deterministic logic. Epistemic Public Infrastructure (EPI) produces classifications, predictions, or inferences through learned parameters.

\begin{table}[htbp]
\centering
\caption{Structural contrast between DPI and EPI}
\label{tab:dpi-epi-contrast}
\begin{tabular}{@{}ll@{}}
\toprule
\textbf{DPI} & \textbf{EPI} \\
\midrule
Deterministic rule execution & Stochastic learned inference \\
Logic transparency sufficient & Weights, data, and representations matter \\
Forkability tied to software & Forkability tied to compute capacity \\
Static certification possible & Performance degrades under distribution shift \\
\bottomrule
\end{tabular}
\end{table}

Corrigibility in EPI therefore requires additional structural conditions beyond those sufficient for deterministic systems.

\paragraph{Tiered Adoption.} EPI systems vary in consequence severity. The framework defines three deployment tiers with increasing compliance requirements:

\begin{table}[htbp]
\centering
\caption{EPI deployment tier classification}
\label{tab:epi-tiers}
\begin{tabular}{@{}lll@{}}
\toprule
\textbf{Tier} & \textbf{Use Cases} & \textbf{Requirements} \\
\midrule
\texttt{trivial} & Translation, summarization & Action boundaries required \\
\texttt{decision\_support} & RAG over government data & Open weights, local data storage \\
\texttt{high\_stakes} & Welfare, policing, judicial & Full LWD-R, human-in-loop \\
\bottomrule
\end{tabular}
\end{table}

The tier classification is a governance decision, not a technical one. Systems must declare their tier; auditors verify appropriateness.

\paragraph{Why ``Epistemic''?} The term is structurally precise. In deterministic DPI, infrastructure manages \textit{what happened}: the \textbf{Ledger State} (a payment was made, an identity was verified). In learned systems, infrastructure manages \textit{what is true}: the \textbf{Inference State} (a person is ``eligible,'' a transaction is ``fraudulent,'' a citizen is a ``risk''). This shift from transactional facts to administrative knowledge constitutes a transfer of \textbf{epistemic sovereignty}.

When a state uses a model to determine eligibility, it is no longer merely processing data; it is producing a \textit{claim to knowledge}. If that claim is hidden in a latent space, the infrastructure has captured the state's ability to know or contest the truth. The ``R'' in the LWD-R Triad (Representation, Section~\ref{sec:code-ai}) is the epistemic anchor: if the categories through which the state ``sees'' the world are proprietary and non-contestable, the state has delegated its interpretive sovereignty. Citizens are not merely locked out of a database; they are locked into a reality they cannot define.

By naming this domain \textbf{Epistemic Public Infrastructure}, the framework asserts that the ``truth'' produced by an AI system must satisfy the same corrigibility tests as a legal fact. The question shifts from ``Does the AI work?'' to ``Is the knowledge produced by this system structurally falsifiable and reversible?''

\paragraph{Architectural Closure.} DPI solves the problem of \textbf{transactions} (Scale and Velocity). EPI solves the problem of \textbf{meaning} (Stochasticity and Representation). Once the infrastructure for creating administrative reality (EPI) is bounded by the same five structural tests as the infrastructure for processing payments (DPI), the framework achieves closure. The theory has scaled from the registry to the world model.

The five structural tests are:

\begin{itemize}[noitemsep]
    \item \textbf{EXIT}: Reversibility of participation
    \item \textbf{CODE}: Inspectability of logic
    \item \textbf{AUDIT}: Independent verification
    \item \textbf{GOVERN}: Binding constraint
    \item \textbf{FORK}: Capacity to reproduce
\end{itemize}

\paragraph{The Cybernetic Foundation.} These tests are not arbitrary normative preferences but derive from control theory. Ashby's Law of Requisite Variety \citep{ashby1956} establishes that a controller must possess at least as many response states as the system it governs. Applied to infrastructure: governance mechanisms must match the variety of harms a system can produce. The five tests operationalize this requirement as a closed feedback loop: EXIT provides the error signal (affected parties can refuse), AUDIT acts as the sensor (errors are detectable), CODE ensures transmission clarity (logic is inspectable), GOVERN serves as the actuator (correction is binding), and FORK provides selection pressure over the entire loop (the system can be reproduced if operators refuse correction).

\paragraph{The Open-Loop Instability Theorem.} When any test fails, the feedback loop opens. Open-loop systems inevitably drift toward instability because errors accumulate without correction. The companion paper proves formally that partial compliance (meeting some tests while failing others) is functionally equivalent to complete failure. A system with EXIT but no GOVERN allows users to detect harm but not correct it; a system with CODE but no FORK allows inspection but not competition. The five tests form an \textit{irreducible} set: removing any one collapses the corrective apparatus.

This paper stress-tests the framework against learned systems: infrastructure where behavior is determined by parameters learned from data (AI/ML) rather than explicit code logic. If the five tests are truly fundamental, they must apply not only to deterministic infrastructure but also to stochastic systems where traditional verification assumptions collapse.

\paragraph{The Challenge.} Learned systems dismantle the assumption that source code is the sole arbiter of behavior. In AI, identical inputs may yield stochastic outputs; behavior is defined by training data; and opacity is often inherent rather than contrived. Yet the collapse of deterministic transparency does not invalidate corrigibility's structural necessity. As systems become less intelligible, external correction becomes \textit{more} critical.

\paragraph{Thesis.} The five corrigibility tests remain invariant for learned systems. What changes is the verification method and the resource barriers to compliance.

\paragraph{Three Structural Pressures.} Learned systems introduce three additional structural pressures on corrigibility: (1) opacity of inference formation, where decision boundaries emerge from training rather than explicit logic; (2) concentration of training resources, where compute asymmetry constrains practical reproduction; and (3) acceleration of automated action, where agentic systems can execute decisions faster than governance mechanisms can respond. The constructs introduced in this paper (LWD-R, Compute Capture, Action Boundary) address each pressure independently.

\paragraph{Scope.} This paper addresses Epistemic Public Infrastructure (Definition~\ref{def:epi}): AI systems deployed at population scale with administrative authority. It does not address consumer AI products where exit is trivial, nor research systems without deployment impact.

\paragraph{Paper Structure.} Section~\ref{sec:invariance} establishes that the five tests remain invariant. Section~\ref{sec:code-ai} extends the CODE test via LWD-R disclosure. Section~\ref{sec:temporal} introduces Variety Drift. Section~\ref{sec:fork-ai} analyzes Compute Capture. Section~\ref{sec:agentic} examines agentic scaling and GDoS. Section~\ref{sec:workflow} addresses workflow sovereignty and WCC. Section~\ref{sec:architecture} provides architecture-specific risk analysis. Section~\ref{sec:discussion} discusses implications and limitations.

\paragraph{Key Contributions.}
\begin{enumerate}[noitemsep]
    \item \textbf{Epistemic Public Infrastructure (EPI)}: Distinguishes AI-mediated administration from deterministic DPI
    \item \textbf{LWD-R}: Four-layer transparency (Logic, Weights, Data, Representation)
    \item \textbf{Ontological Capture}: When a model's categories become non-contestable facts
    \item \textbf{Variety Drift}: Why AI corrigibility is a persistence condition
    \item \textbf{Compute Capture}: When ``open weights'' fails FORK
    \item \textbf{GDoS}: Governance collapse under agentic scaling
    \item \textbf{WCC}: Quantifies epistemic delegation risk
    \item \textbf{Action Boundary}: Deterministic envelope around stochastic inference
\end{enumerate}

%% ============================================
%% SECTION 2: INVARIANCE
%% ============================================

\section{The Invariance of the Standard}
\label{sec:invariance}

\begin{proposition}
The five corrigibility tests remain invariant for learned systems. What changes is the verification method.
\end{proposition}

The structural requirement for corrigibility derives from Ashby's Law of Requisite Variety: a controller must have at least as many response states as the system it governs \citep{ashby1956}. This requirement is independent of whether the controlled system is deterministic or stochastic. An AI model that affects citizen outcomes at population scale generates the same variety of potential harms as deterministic infrastructure; therefore, the governance apparatus must possess equivalent corrective variety.

The table below illustrates how the mechanical verification of each test adapts to the probabilistic nature of AI without softening the structural requirement.

\begin{table}[htbp]
\centering
\small
\begin{tabular}{@{}lp{3.5cm}p{4.5cm}p{3cm}@{}}
\toprule
\textbf{Test} & \textbf{DPI Verification} & \textbf{EPI Verification} & \textbf{Key Addition} \\
\midrule
EXIT & Can refuse the system & Disclosure + non-AI alternative; appeal to human guaranteed & Human fallback \\
CODE & Source code determines behavior & LWD-R: Logic, Weights, Data, Representation layers disclosed & \textbf{R} (ontology) \\
AUDIT & Inspect logic and outputs & Statistical bounds + continuous monitoring; variety drift measurement & Drift tracking \\
GOVERN & Maintainer and RFC process & Governance over objectives, value tradeoffs, and category definitions & Ontology governance \\
FORK & Copy code and deploy & Resource forkability: compute + data + training pipeline & Compute access \\
\bottomrule
\end{tabular}
\caption{Verification methods for DPI (deterministic) vs. EPI (learned) systems. The standard remains invariant; verification adapts.}
\label{tab:verification}
\end{table}

\paragraph{Key Insight.} The shift from deterministic to learned systems does not weaken the tests; it \textit{strengthens} the case for structural corrigibility. When behavior becomes less predictable, the capacity for external correction becomes more essential. A deterministic system that fails silently can be debugged; a learned system that fails silently accumulates error without detection.

%% ============================================
%% SECTION 3: CODE FOR AI (LWD-R)
%% ============================================

\section{Transparency Failure in Learned Systems}
\label{sec:code-ai}

In deterministic infrastructure, disclosure of operative logic is sufficient to enable inspection. In learned systems, logic is distributed across architecture, parameters, training data, and representational schemas. Disclosure of source code alone does not reveal the operative decision boundary.

\begin{principle}[Source Distribution]
Publishing inference code alone does not satisfy the CODE test. It allows an auditor to run the machine, but not to understand why it produces specific outputs.
\end{principle}

To satisfy the CODE condition in learned systems, four components must be disclosed:

\begin{itemize}[noitemsep]
    \item \textbf{Logic (L)}: Model architecture and inference procedures
    \item \textbf{Weights (W)}: Learned parameters
    \item \textbf{Data (D)}: Training datasets with provenance
    \item \textbf{Representation (R)}: Ontological categories and label schemas
\end{itemize}

We refer to this as \textbf{LWD-R transparency}. Omission of any component prevents meaningful external reconstruction of epistemic behavior.

\paragraph{Scope and Privacy.} LWD-R defines transparency sufficient for reconstruction of epistemic behavior. It does not require public disclosure of sensitive data in raw form, but requires verifiable disclosure sufficient for independent audit and reproduction. Data manifests with provenance, statistical summaries, and representative samples may satisfy the D requirement without exposing individual records. The standard is auditability, not publication.

To bridge this gap, compliance requires structured disclosure across the training pipeline, such as \textbf{Data Sheets for Datasets} \citep{gebru2021} and \textbf{Model Cards} \citep{mitchell2019}. These artifacts function as the ``source code'' for structural evaluation; withholding them renders the system a black box.

\subsection{The LWD-R Disclosure Requirement}

For world models and agentic systems, the traditional transparency requirements (code, weights, data) are necessary but insufficient. Consider what it means to ``inspect'' a welfare eligibility model: seeing the inference code tells you how predictions are computed, but not \textit{why} certain applicants are classified as high-risk. That ``why'' is encoded across four distinct layers.

\begin{definition}[LWD-R: The Four Layers of AI Transparency]
\label{def:lwdr}
A learned system satisfies LWD-R disclosure if the following artifacts are publicly available:
\begin{itemize}[noitemsep]
    \item \textbf{L (Logic)}: Model architecture and inference code (\textit{how} predictions are computed)
    \item \textbf{W (Weights)}: Trained parameters (the learned patterns that produce outputs)
    \item \textbf{D (Data)}: Training corpus with provenance (\textit{what} the model learned from)
    \item \textbf{R (Representation)}: The categorical schema (\textit{how the model sees the world})
\end{itemize}
\end{definition}

\paragraph{Why Four Layers?} The first three (L, W, D) are familiar from ``open source AI'' discourse. The fourth, \textbf{Representation (R)}, is this framework's principal contribution to AI transparency discourse.

\paragraph{Why ``R'' Is the Novel Contribution.} Existing ``open AI'' frameworks focus on L, W, and D. The Representation layer addresses a distinct failure mode: a model may publish weights, code, and data manifests while encoding \textit{non-contestable categories} in its latent space. When an eligibility model classifies citizens as ``high-risk,'' that category becomes an administrative fact. If the category itself cannot be challenged, the model has captured the state's interpretive sovereignty.

\textbf{R} is not transparency about \textit{how} the model works; it is transparency about \textit{how the model sees the world}. These categories (the model's ontology) are often hidden in the latent space, undocumented and non-contestable. Failure to disclose \textbf{R} means the model's worldview is structurally opaque, even if weights are published.

\subsection{Ontological Capture}
\label{subsec:rcr}

When a society becomes dependent on a model's categorical schema without structural means to contest those categories, we call this \textbf{Ontological Capture}.

\paragraph{The Problem.} A welfare eligibility model classifies applicants into ``low-risk'' and ``high-risk.'' A fraud detection system labels transactions as ``suspicious'' or ``normal.'' A world model parses geopolitical actors into ``threat'' categories. These classifications are not neutral descriptions; they are \textit{definitional acts} that shape administrative reality. If the categories themselves cannot be challenged (if no one outside the vendor can ask ``why is this person high-risk?'' and receive a contestable answer), the state has surrendered its interpretive sovereignty.

\begin{definition}[Ontological Capture]
\label{def:ontological-capture}
Let $\Delta O$ represent the \textbf{ontological contestability} of a system: the degree to which affected parties can challenge, modify, or replace the representation schema. Ontological Capture occurs when $\Delta O \to 0$: when the model's categories become non-contestable administrative facts.
\end{definition}

\paragraph{Connection to GOVERN.} Ontological capture occurs when affected participants cannot contest or modify the representational categories through which they are classified. This is a failure of the GOVERN layer within learned systems: the categories themselves (``high-risk,'' ``eligible,'' ``fraudulent'') become binding without participatory constraint. GOVERN in EPI requires governance not only over system behavior but over the ontological schema that defines the system's view of reality.

\paragraph{Implication for CODE.} Under Ontological Capture, the system fails CODE not because the weights are hidden, but because the \textit{meaning structure} is opaque. Compliance requires not merely weight publication but \textbf{representation schema disclosure}: machine-readable specifications of the categories the model employs, with documented justification for category boundaries.

%% ============================================
%% SECTION 4: TEMPORAL VARIETY GAP
%% ============================================

\section{Variety Drift: Why AI Corrigibility Is a Persistence Condition}
\label{sec:temporal}

Traditional software remains static until explicitly patched. A tax calculation program written in 2020 will compute the same result in 2030, unless someone changes the code. Learned systems are different: they are \textit{frozen snapshots} of a world that keeps moving.

\paragraph{The Problem.} An AI model trained on 2024 data encodes 2024 patterns. By 2026, the world has changed: new fraud techniques emerge, economic conditions shift, language evolves. The model's ``variety'' (its capacity to handle diverse situations) was fixed at training. The world's variety keeps expanding. The gap between them grows over time.

\begin{definition}[Variety Drift]
\label{def:variety-drift}
Let $V_{\text{world}}(t)$ be the variety of situations the environment presents at time $t$, and $V_{\text{model}}(\theta)$ be the variety the model can handle, fixed at training. \textbf{Variety Drift} is the growing gap:
\begin{equation}
\Delta(t) = V_{\text{world}}(t) - V_{\text{model}}(\theta)
\label{eq:variety-gap}
\end{equation}
\end{definition}

\paragraph{Why This Matters.} Without mechanisms for retraining accessible to the community, variety drift widens indefinitely. A model that passed all five corrigibility tests at deployment may fail them two years later. Not because anyone changed anything, but because the world moved and the model did not. \textbf{AI corrigibility is not a deployment certification; it is a persistence condition.}

\paragraph{Scope.} Variety drift represents temporal degradation of epistemic adequacy under distributional shift. It is not unique to learned systems (deterministic systems also face changing contexts), but is amplified by scale and opacity. In learned systems, the gap between model and reality is invisible without active monitoring; in deterministic systems, logic mismatch can at least be traced to specific rules.

\begin{figure}[htbp]
\centering
\begin{tikzpicture}[x=0.8cm, y=0.5cm, font=\sffamily]

    % Axes
    \draw[->, thick, gray] (0,0) -- (10,0) node[right] {Time ($t$)};
    \draw[->, thick, gray] (0,0) -- (0,8) node[above] {Variety ($V$)};

    % Real World Complexity (The Disturbance) - Curving up
    \draw[thick, forkcolor, smooth] plot coordinates {(0,1) (2,1.8) (4,3.5) (6,4.5) (8,7.5) (9.5, 9)};
    \node[forkcolor, anchor=west, font=\scriptsize] at (9.5, 9) {$V(\text{Environment})$};

    % AI Model State (The Controller) - Step Function
    \draw[thick, codecolor] (0,1) -- (4,1);
    \draw[dashed, gray] (4,1) -- (4,3.5);
    \draw[thick, codecolor] (4,3.5) -- (8,3.5);
    \draw[dashed, gray] (8,3.5) -- (8,7.5);
    \draw[thick, codecolor] (8,7.5) -- (9.5,7.5);

    % Labels for Model
    \node[codecolor, anchor=north, font=\scriptsize] at (2,1) {Model $\theta_1$};
    \node[codecolor, anchor=north, font=\scriptsize] at (6,3.5) {Model $\theta_2$};

    % The Gap (Error)
    \fill[red!10] (0,1) -- (4,1) -- (4,3.5) -- plot[smooth] coordinates {(4,3.5) (2,1.8) (0,1)} -- cycle;
    \fill[red!20] (4,3.5) -- (8,3.5) -- (8,7.5) -- plot[smooth] coordinates {(8,7.5) (6,4.5) (4,3.5)} -- cycle;

    % Annotations
    \node[anchor=west, text=red!80!black, font=\scriptsize] at (0.5, 2.5) {$\Delta(t)$};
    \draw[->, red!80!black] (1.2, 2.3) -- (1.5, 1.5);

    \node[align=center, font=\tiny, gray] at (4, -1) {\textbf{Retraining}\\(FORK)};
    \draw[->, gray, dotted] (4, -0.3) -- (4, 0.8);

\end{tikzpicture}
\caption{\textbf{The Temporal Variety Gap.} Environmental variety $V(e)$ evolves continuously (curve). A fixed model $\theta$ (stepped line) diverges from reality. The shaded area represents accumulated error. Only \textbf{FORK} rights allow the gap to be closed via retraining.}
\label{fig:variety_gap}
\end{figure}

\begin{claim}
A learned system that passes all five corrigibility tests at deployment may fail them over time if the community lacks the resources to retrain. Corrigibility for AI is a dynamic property, not a static certification.
\end{claim}

%% ============================================
%% SECTION 5: FORK AND COMPUTE CAPTURE
%% ============================================

\section{Resource Barriers to Forkability}
\label{sec:fork-ai}

In traditional software, forking is cheap: copy the code, set up a server, deploy. The cost is measured in storage and labor. In AI, forking is expensive: the cost is measured in \textit{compute}, the GPU-hours required to retrain a model from scratch.

\paragraph{The Structural Difference.} When Elastic changed its license, AWS forked Elasticsearch and the community continued development. This was possible because the resource barrier was low. When Meta releases Llama weights, can the community ``fork'' Llama in the same sense? They can \textit{run} it (inference). They can \textit{adapt} it (fine-tuning). But can they \textit{reproduce} it from scratch if Meta's training data turns out to be problematic? For frontier models, the answer is no: the compute cost exceeds what any non-corporate actor can access.

\begin{table}[htbp]
\centering
\begin{tabular}{@{}lll@{}}
\toprule
\textbf{Layer} & \textbf{Cost} & \textbf{Barrier} \\
\midrule
Inference code & Minimal & None \\
Running open weights & Low & Minor \\
Fine-tuning & Moderate & Economic \\
Retraining from scratch & Very high & Structural \\
\bottomrule
\end{tabular}
\caption{Resource barriers to forking learned systems}
\label{tab:fork-barriers}
\end{table}

\begin{claim}
Legal permission to fork without access to compute or data is meaningless. This is an infrastructure question, not a licensing question.
\end{claim}

\subsection{Compute Capture}
\label{subsec:compute-capture}

\paragraph{The Problem.} A company releases model weights under an open license. Headlines announce ``open source AI.'' But the training run cost \$100 million in compute. No university, no government research lab, no coalition of independent researchers can afford to retrain it. If the model exhibits problematic behavior traceable to training data, the community can observe the symptoms but cannot fix the cause. They can run the model; they cannot reproduce it.

\begin{definition}[Compute Capture]
\label{def:compute-capture}
\textbf{Compute Capture} occurs when the compute cost to retrain a model ($C_{\text{train}}$) exceeds the compute accessible to non-operator actors ($C_{\text{accessible}}$) by more than a policy-determined threshold:
\begin{equation}
C_{\text{train}} > \kappa \cdot C_{\text{accessible}}
\label{eq:compute-capture}
\end{equation}
where $\kappa$ is an accessibility multiplier (typically $1 \leq \kappa \leq 2$). Under Compute Capture, releasing weights provides \textbf{inference forkability} (you can run it) without \textbf{training forkability} (you can reproduce it).
\end{definition}

\paragraph{Structural Claim.} Compute capture does not imply that large-scale training is illegitimate. It asserts that when reproduction exceeds accessible resource bounds, forkability becomes structurally constrained, weakening the replacement layer. The claim is architectural: legal permission to fork is meaningless without physical capacity to fork.

\paragraph{Why This Matters.} Legal permission to fork is meaningless without physical capacity to fork. If training cost exceeds 100$\times$ global academic compute capacity, the FORK test fails regardless of licensing terms. This is an infrastructure question, not a licensing question.

The formal FORK determination for training becomes:
\begin{equation}
\text{FORK}_{\text{training}} = \mathbf{1}[C_{\text{train}} \leq \kappa \cdot C_{\text{accessible}}] \land \mathbf{1}[\text{data\_manifest\_available}]
\label{eq:fork-training}
\end{equation}
where $\mathbf{1}[\cdot]$ is the indicator function.

\subsection{Compute as Governance Infrastructure}

This analysis implies that \textbf{compute is itself DPI}. Legal forkability is insufficient absent practical compute accessibility. Architectural responses include:
\begin{enumerate}[noitemsep]
    \item \textbf{Public Compute Nodes}: State-sponsored clusters with multi-stakeholder governance
    \item \textbf{Mandatory Cost Disclosure}: FLOP estimates and data manifests
    \item \textbf{Federated Training Pools}: Distributed training across independent nodes
    \item \textbf{Compute Vouchers}: Grants for reproduction attempts
\end{enumerate}

\subsection{Minimum Evidence for Learned System Forkability}
\label{subsec:fork-evidence}

To prevent ``FORK fails'' from becoming a rhetorical claim rather than a verifiable determination, learned systems must satisfy specific evidentiary requirements:

\begin{table}[htbp]
\centering
\small
\caption{Evidence requirements for FORK in learned systems}
\label{tab:ai-fork-evidence}
\begin{tabular}{@{}llp{5.5cm}@{}}
\toprule
\textbf{Artifact} & \textbf{Status} & \textbf{Verification} \\
\midrule
Inference code & Required & Open source license; build reproducibility \\
Model weights & Required & Public download; hash verification \\
Training data manifest & Required & Data Sheets \citep{gebru2021} with provenance \\
Training checkpoint & Recommended & Enables fine-tuning without full retrain \\
Training code & Required & Includes preprocessing, hyperparameters \\
Compute cost estimate & Required & FLOP estimate or cloud cost documentation \\
Successful reproduction & Decisive & Third-party retrain within 10\% of benchmark \\
\bottomrule
\end{tabular}
\end{table}

\paragraph{Forkability Determination.} A learned system passes FORK if:
\begin{enumerate}[noitemsep]
    \item All ``Required'' artifacts are publicly available under permissive terms
    \item The compute cost estimate demonstrates feasibility for at least one non-operator entity
    \item Either a successful third-party reproduction exists, OR the training code + data manifest enables reproduction in principle
\end{enumerate}

\subsection{Inference vs. Training Forkability}

Systems that publish weights but withhold training data or code achieve only \textbf{inference forkability}, not \textbf{training forkability}. Inference forkability permits running the model; training forkability permits correcting it. Only the latter satisfies FORK for governance purposes.

\paragraph{Concrete Examples.} As of 2026, most ``open'' models achieve only inference forkability:
\begin{itemize}[noitemsep]
    \item \textbf{Inference forkability only:} Llama 3 (Meta), DeepSeek, Gemma (Google), Mistral. Weights are published, but training data remains proprietary and full training code is not disclosed. Users can run and fine-tune these models but cannot independently reproduce or audit the training process.
    \item \textbf{Training forkability:} OLMo (Allen Institute), Pythia (EleutherAI). These projects publish weights, complete training code, and full data manifests (Dolma and The Pile respectively). Third parties have successfully reproduced training runs, satisfying the FORK test.
\end{itemize}

The distinction matters: when Llama or Gemma exhibit unexpected behavior, the community can observe symptoms but cannot diagnose root causes in training data. When OLMo exhibits unexpected behavior, researchers can trace the issue to specific training examples. \textbf{Open weights with closed training pipelines is inference transparency without correction capacity.}

\paragraph{Alignment with OSI Definition.} The Open Source Initiative's Open Source AI Definition 1.0 \citep{osi2024} codifies this distinction, requiring: (1) data information sufficient for recreation, (2) complete training code, and (3) model parameters. By this standard, Llama, Gemma, DeepSeek, and Mistral do not qualify as open source AI despite marketing claims. The EU AI Act \citep{euaiact2024} similarly distinguishes between models with varying transparency obligations. This framework's FORK test operationalizes the same principle: without training forkability, governance correction is structurally impossible.

%% ============================================
%% SECTION 6: AGENTIC SCALING
%% ============================================

\section{Agentic Scaling and Governance Denial of Service}
\label{sec:agentic}

The preceding sections establish that learned systems face unique barriers to corrigibility: opaque training pipelines (violating CODE), frozen variety (the Temporal Variety Gap), and compute-gated reproduction (violating FORK). These challenges intensify as AI moves from isolated models to autonomous agents operating within infrastructure. As DPI transitions to support the ``Agentic Economy,'' corrigibility moves from a moral preference to an engineering requirement. Autonomous agents lack the biological resilience to handle bureaucracy.

\begin{principle}[Structural vs. Alignment Distinction]
This framework distinguishes structural corrigibility from ``AI Alignment.'' The AI safety literature has extensively studied alignment (ensuring systems pursue operator-intended goals) \citep{soares2017} and the ``off-switch'' problem \citep{hadfield2017}. These framings center on \textit{operator control}. Alignment asks: \textit{Can the operator control the system?} Corrigibility asks: \textit{Can the affected subject correct the system?} Structural legitimacy requires the latter.
\end{principle}

\subsection{Corrigibility as an API Contract for Agents}

\begin{itemize}[noitemsep]
    \item \textbf{EXIT is Exception Handling}: To an agent, a mandatory system with no exit path appears as a logic deadlock.
    \item \textbf{CODE is Documentation}: Agents require inspectability of logic to form valid plans.
    \item \textbf{AUDIT is Observability}: An optimization agent cannot function without a reliable failure signal.
    \item \textbf{GOVERN is Constraint Discovery}: Advisory guidelines result in undefined behavior for software agents.
\end{itemize}

\begin{claim}
Incorrigible infrastructure is structurally incompatible with AI automation.
\end{claim}

\subsection{The Human Friction Buffer}

Incorrigible infrastructures may appear stable under human load because humans absorb friction through delay, compliance, and informal workaround. Autonomous agents remove this buffering layer.

\paragraph{Illustrative Scenario.} Consider an autonomous procurement agent that interacts with a mandatory payment API. On encountering a compliance flag, the agent receives an opaque error code. No machine-readable error taxonomy exists, and dispute resolution is human-mediated. The agent retries, triggering repeated failures and queue saturation. Human agents absorb such friction via manual override and phone calls; fleets of autonomous agents cannot. The result is a governance denial-of-service (GDoS) that cascades across the payment rail.

\subsection{Governance Denial of Service (GDoS)}

\paragraph{The Problem.} A human user encounters a bureaucratic error and waits on hold for 45 minutes. Frustrating, but the system survives. Now imagine 10,000 autonomous agents encountering the same error simultaneously. They don't wait patiently; they retry. They don't call a human; they escalate programmatically. The governance mechanism, designed for human patience, collapses under machine speed.

\begin{definition}[Governance Denial of Service]
\label{def:gdos}
\textbf{GDoS} occurs when the volume, speed, and strategic diversity of agent actions exceeds the system's capacity to audit, interpret, and correct behavior. It is not a security attack but a control-theoretic inevitability: when error accumulation rate ($\frac{dE}{dt}$) exceeds correction rate ($\frac{dC}{dt}$), the system diverges.
\end{definition}

\paragraph{Why This Is Inevitable.} Agents accelerate error accumulation; centralized governance cannot accelerate correction proportionally. Systems lacking distributed audit and rapid, binding governance lack the requisite variety to remain stable under agentic scaling.

\begin{claim}
Agentic scaling amplifies the correction velocity inequality. Systems that are marginally stable under human interaction rates become catastrophically unstable under agent interaction rates. The human friction buffer is not a feature; it is a hidden subsidy that masks architectural incorrigibility.
\end{claim}

\subsection{Agentic Interoperability Requirements}

For an autonomous agent $\mathcal{A}$ to operate reliably within infrastructure $\mathcal{I}$, the following predicates must be strictly True:

\begin{enumerate}[noitemsep]
    \item \textbf{Reversibility:} $\forall$ state $s_t \in \mathcal{S}$, $\exists$ transition $s_{t+1}$ (via \texttt{EXIT}) such that $\text{Cost}(s_{t+1}) < \infty$. (No deadlocks).
    \item \textbf{Decidability:} $\forall$ input $x$, $\text{Compute}(x) \to \{0,1\}$ is inspectable via \textbf{CODE}. (No hidden constraints).
\end{enumerate}

Failure of these predicates renders $\mathcal{I}$ a ``Hazmat Environment'' for automated agents, effectively blocking the agentic economy.

\begin{figure}[htbp]
\centering
\begin{tikzpicture}[node distance=2cm, auto, >=latex, thick]
    \tikzstyle{state} = [draw, ellipse, minimum width=1.5cm, minimum height=1cm]
    \tikzstyle{blocked} = [draw, ellipse, minimum width=1.5cm, minimum height=1cm, fill=red!20]

    \node [state] (idle) {Idle};
    \node [state, right of=idle, node distance=3cm] (call) {Call API};
    \node [state, right of=call, node distance=3cm] (success) {Success};
    \node [blocked, below of=call, node distance=2cm] (blocked) {Blocked};
    \node [state, below of=success, node distance=2cm] (error) {Error};

    \draw [->] (idle) -- node {\scriptsize request} (call);
    \draw [->] (call) -- node {\scriptsize 2xx} (success);
    \draw [->] (success) -- node {\scriptsize done} ++(1.5,0);
    \draw [->] (call) -- node[left] {\scriptsize blocked} (blocked);
    \draw [->] (call) to[loop above, looseness=5] node {\scriptsize retry} (call);
    \draw [->] (call) -- node {\scriptsize timeout} (error);
    \draw [->] (blocked) -- node[below] {\scriptsize (No Exit)} ++(-2,0);
\end{tikzpicture}
\caption{\textbf{Agentic automaton.} Without EXIT and observable error codes, agent loops indefinitely, times out, or escalates incorrectly.}
\label{fig:agentic-automaton}
\end{figure}

\subsection{The Action Boundary Protocol}
\label{subsec:action-boundary}

As infrastructure transitions to the Rule of the Workflow, the fundamental unit of governance can no longer be the static database record, nor can it be the stochastic model weight. To subject autonomous systems to democratic constraint, the architecture must adopt a new structural primitive: the \textbf{Action Boundary}.

\begin{principle}[Action Boundary Protocol]
If the state cannot govern the neural network's weights, it must govern the \textbf{action boundary}: the deterministic envelope that wraps stochastic inference. The Action Boundary Protocol specifies: (1) hard-coded preconditions, (2) deterministic tool execution, and (3) explicit termination criteria to prevent Governance Denial of Service (GDoS).
\end{principle}

\paragraph{Architectural Definition.} The Action Boundary Protocol separates probabilistic inference from deterministic authorization. Inference may generate proposals; execution must remain governed by policy constraints independent of inference confidence. This preserves constraint-layer corrigibility under agentic scaling.

In agentic systems, learned inference should not directly trigger irreversible action. Instead, inference outputs must pass through a deterministic validation layer that enforces policy constraints and safety rules.

\begin{figure}[htbp]
\centering
\begin{tikzpicture}[
    node distance=0.9cm,
    box/.style={rectangle, draw, rounded corners, minimum height=1cm, minimum width=2.4cm, font=\small\sffamily, align=center},
    stochastic/.style={box, fill=orange!15, draw=orange!60},
    deterministic/.style={box, fill=green!15, draw=green!60},
    reject/.style={box, fill=red!15, draw=red!60},
    arrow/.style={-{Stealth[length=2.5mm]}, thick}
]

% Main flow
\node[stochastic] (llm) {LLM\\Inference};
\node[box, right=of llm] (action) {Action\\Proposal};
\node[deterministic, right=of action] (validator) {Deterministic\\Validator};
\node[deterministic, right=of validator] (exec) {Execution};

% Reject path
\node[reject, below=1.2cm of validator] (reject) {REJECT\\(Logged)};

% Arrows - main flow
\draw[arrow] (llm) -- (action);
\draw[arrow] (action) -- (validator);
\draw[arrow] (validator) -- (exec);
\draw[arrow, red!60] (validator) -- (reject);

% Feedback arrow from reject to llm
\draw[arrow, dashed, blue!60] (reject.west) -- ++(-3.5,0) |- (llm.south);
\node[font=\tiny\sffamily, text=blue!60, below=0.1cm of llm] {Learning feedback};

% Domain labels
\node[above=0.4cm of llm, font=\scriptsize\sffamily\bfseries, text=orange!70!black] {INFERENCE DOMAIN};
\node[above=0.4cm of exec, font=\scriptsize\sffamily\bfseries, text=green!50!black] {EXECUTION DOMAIN};

% Action Boundary - vertical separator (visually strong)
\draw[line width=2pt, blue!80] ($(action.east)+(0.45,1.4)$) -- ($(action.east)+(0.45,-2.2)$);
\node[font=\scriptsize\sffamily\bfseries, text=blue!80, rotate=90, anchor=south] at ($(action.east)+(0.75,-0.4)$) {ACTION BOUNDARY};

% GOVERN label and deterministic policy note
\node[below=0.1cm of validator, font=\tiny\sffamily, text=green!50!black] {(GOVERN)};
\node[above=0.15cm of validator, font=\tiny\sffamily\itshape, text=green!50!black] {Deterministic policy layer};

\end{tikzpicture}
\caption{\textbf{Action Boundary Protocol.} Probabilistic inference (orange) produces action proposals. The \textbf{deterministic policy layer} (green, implementing GOVERN) enforces hard constraints before execution. Invalid actions are rejected, logged, and fed back for learning. The vertical \textbf{Action Boundary} separates the inference domain from the execution domain: the architectural instantiation of corrigibility in agentic systems.}
\label{fig:action-boundary}
\end{figure}

The validator enforces hard constraints independent of probabilistic inference. Corrigibility in agentic systems therefore requires separation between epistemic generation and action authorization.

The engineering instantiation of this protocol already exists: the \textbf{Agentic Skills} standard \citep{agentskills2025}.

\paragraph{Agentic Skills (agentskills.io).} The open \texttt{SKILL.md} specification defines a bounded, machine-readable orchestration protocol that encapsulates stochastic inference within deterministic action boundaries. Adopted by 18+ AI agent platforms (OpenAI Codex, GitHub Copilot, Cursor, and others), this standard provides exactly the governance primitive required for the Action Boundary Protocol. We adopt it here as the technical instantiation of our theoretical framework.

\paragraph{The SBOM Analogy.} For security-minded readers: Agentic Skills function like a \textbf{Software Bill of Materials (SBOM)} for AI actions. Just as SBOMs enumerate software dependencies to enable vulnerability auditing, skill manifests enumerate action boundaries to enable governance auditing. The same transparency that enables correction also creates attack surface. This tradeoff is inherent to all deterministic governance systems (see Limitations, Section~\ref{sec:discussion}).

A skill is a directory containing a \texttt{SKILL.md} file with YAML frontmatter and Markdown instructions:

\begin{verbatim}
skill-name/
+-- SKILL.md          # Required: frontmatter + instructions
+-- scripts/          # Optional: executable code
+-- references/       # Optional: additional documentation
+-- assets/           # Optional: templates, schemas
\end{verbatim}

\paragraph{Frontmatter as Corrigibility Contract.} The YAML frontmatter encodes the governance contract:

\begin{verbatim}
---
name: welfare-triage           # Required, 1-64 chars
description: Determines eligibility for welfare programs.
             Use when processing benefit applications.
license: Apache-2.0
compatibility: Requires access to citizen registry API
allowed-tools: Bash(curl:*) Read Write  # GOVERN constraint
metadata:
  author: state-welfare-dept
  version: "2.1"
---
\end{verbatim}

The \texttt{allowed-tools} field is the \textbf{cybernetic anchor}: it specifies which tools the agent may invoke, hard-coded outside the LLM's context window. An LLM cannot prompt-inject past this constraint. It is enforced by the execution environment, not by the model's alignment.

\paragraph{Progressive Disclosure Architecture.} Skills use progressive disclosure to manage context efficiently, directly implementing the LWD-R requirement:

\begin{table}[htbp]
\centering
\small
\caption{SKILL.md progressive disclosure layers}
\label{tab:skill-disclosure}
\begin{tabular}{@{}llll@{}}
\toprule
\textbf{Layer} & \textbf{Tokens} & \textbf{When Loaded} & \textbf{LWD-R Mapping} \\
\midrule
Metadata & $\sim$100 & Startup (all skills) & L (Logic discovery) \\
Instructions & $<$5000 & Activation & R (Representation) \\
Resources & As needed & Execution & D (Data provenance) \\
\bottomrule
\end{tabular}
\end{table}

\subsubsection{Bounding Stochasticity}

In learned systems, behavior is probabilistic. However, public infrastructure requires deterministic accountability. The Agentic Skills architecture resolves this tension by separating the \textit{reasoning engine} (the stochastic model) from the \textit{action boundary} (the deterministic skill definition).

\begin{principle}[Action Boundary Governance]
If the state cannot govern the neural network's weights, it must govern the action boundary. Agentic Skills provide the deterministic envelope that wraps around probabilistic inference.
\end{principle}

When an agent attempts a state transition, the workflow orchestration does not rely on the model's internal alignment. Instead, it invokes a defined Skill that imposes exogenous constraints:

\begin{enumerate}[noitemsep]
    \item \textbf{Exception Handling (EXIT)}: The skill definition dictates explicit fallback paths and termination triggers, preventing the infinite loops that cause Governance Denial of Service (GDoS).
    \item \textbf{Ontological Disclosure (CODE)}: By encoding the agent's workflow into declarative, human-readable files (e.g., YAML/Markdown skill definitions), the system's operational ontology becomes public. This satisfies the Representation (R) requirement of the LWD-R triad.
    \item \textbf{Telemetry (AUDIT)}: Skills enforce structured inputs and outputs, transforming opaque ``conversations'' into discrete, auditable transaction logs.
    \item \textbf{Binding Rules (GOVERN)}: The execution of the skill acts as the non-overrideable controller. An LLM cannot prompt-inject its way out of a hard-coded skill constraint.
\end{enumerate}

\subsubsection{Defeating Workflow Capture via Open Skills}

The critical danger of agentic DPI is Epistemic Delegation: the loss of state sovereignty to proprietary orchestration layers. We previously defined the Workflow Capture Coefficient (WCC) as a function of reproduction cost, ontological substitutability, and portability.

Standardized Agentic Skills fundamentally disrupt this capture vector. If a state defines its administrative processes using open skill protocols rather than vendor-locked pipelines:

\begin{itemize}[noitemsep]
    \item $C_{\text{workflow}} \to 0$: The orchestration logic is transparent and easily replicable
    \item $O_{\text{substitutability}} \to 1$: Administrative categories are defined in open skill schemas, not hidden in a vendor's latent space
    \item $M_{\text{portability}} \to 1$: The state can swap the underlying LLM without losing or breaking the workflow logic
\end{itemize}

Consequently, adopting an open Agentic Skills architecture drives the WCC toward zero, preserving the state's interpretive sovereignty and satisfying the \textbf{FORK} test for the agentic era.

\begin{claim}
You don't govern the AI; you govern the Skills the AI is allowed to use. This is the architectural translation of corrigibility into the agentic domain.
\end{claim}

\paragraph{Critical Limitation: Action Boundary Is Containment, Not Cure.} The Action Boundary Protocol is a \textit{necessary but not sufficient} condition for EPI corrigibility. It bounds the ``act'' but does not resolve deeper structural deficits:

\begin{enumerate}[noitemsep]
    \item \textbf{Inference-layer opacity persists.} The model's internal representation (the categories it uses to \textit{interpret} inputs before selecting a skill) remains hidden. You govern what the model \textit{does}, not what it \textit{thinks}. Representation Capture Risk (RCR) operates upstream of the Action Boundary.

    \item \textbf{Compute capture blocks training forkability.} Skill portability enables swapping orchestration logic, but the underlying model remains compute-gated. If the LLM itself cannot be retrained by non-operators, the Action Boundary governs a black box.

    \item \textbf{Temporal Variety Gap is unaddressed.} The model drifts from reality between retraining cycles. Bounding actions does not prevent the accumulating mismatch between frozen model variety and evolving environmental variety.

    \item \textbf{Skill definitions encode ontological choices.} The categories embedded in skill schemas (``eligible,'' ``fraudulent,'' ``risk'') are themselves epistemic commitments. Open skill protocols shift capture from the LLM to the skill author; they do not eliminate it.
\end{enumerate}

The Action Boundary is therefore a \textbf{governance floor}, not a governance ceiling. It prevents the worst failure modes (GDoS, prompt injection, unbounded action) while leaving the deeper problem of epistemic sovereignty partially unresolved. Full EPI corrigibility requires the Action Boundary \textit{plus} LWD-R disclosure \textit{plus} training forkability: the complete stack, not any single layer.

%% ============================================
%% SECTION 7: WORKFLOW SOVEREIGNTY
%% ============================================

\section{Workflow Sovereignty and the Rule of Workflow}
\label{sec:workflow}

The final mutation of digital infrastructure is the transition from static registries to \textbf{Workflow Orchestration layers} (e.g., Palantir-style data fusion or agentic AI pipelines). In this regime, the operating layer of the state is no longer a ledger of records, but an orchestration of decisions.

\subsection{From Rule of Law to Rule of Workflow}

In classical governance, policy is defined in statute, executed by bureaucracy, and reviewed by courts. In a ``Rule of Workflow'' regime, policy becomes emergent from the orchestration logic: data schemas define reality, prompts define interpretation, and API contracts define valid inputs.

\textbf{Workflow Capture} is formalized as follows. Let $W$ be the orchestration layer, $M$ the underlying model, $S$ the state policy function, and $D$ the decision output. While classical governance assumes $S \rightarrow D$, AI-native governance operates as:
\begin{equation}
S \rightarrow W(M, \text{Data}) \rightarrow D
\end{equation}

If $W$ is proprietary and non-forkable, the state's ability to alter policy ($\Delta S$) vanishes because outcomes are constrained by the workflow's technical affordances. Policy debate is reduced to parameter tuning within a vendor-controlled ontology.

\begin{figure}[htbp]
\centering
\begin{tikzpicture}[scale=0.9, font=\sffamily\small,
    block/.style={rectangle, rounded corners, draw=gray!70, fill=gray!10, minimum width=1.8cm, minimum height=0.8cm},
    wblock/.style={rectangle, rounded corners, draw=forkcolor, fill=forkcolor!15, minimum width=2.5cm, minimum height=1cm, font=\bfseries},
    arrow/.style={-{Stealth[length=2mm]}, thick}
]

% Classical Governance (top)
\node[anchor=west, font=\bfseries] at (-4, 2.5) {Classical:};
\node[block, fill=governcolor!20] (S1) at (0, 2.5) {Policy $S$};
\node[block] (D1) at (4, 2.5) {Decision $D$};
\draw[arrow, governcolor] (S1) -- (D1) node[midway, above, font=\scriptsize] {direct};

% AI-Native Governance (bottom)
\node[anchor=west, font=\bfseries] at (-4, 0) {AI-Native:};
\node[block, fill=governcolor!20] (S2) at (-1, 0) {Policy $S$};
\node[wblock] (W) at (2.5, 0) {Workflow $W$};
\node[block] (D2) at (6, 0) {Decision $D$};

% Internal to W
\node[font=\tiny, text=gray] at (2.5, -0.6) {Model $M$ + Data};

\draw[arrow, governcolor] (S2) -- (W);
\draw[arrow, forkcolor] (W) -- (D2);

% Capture indicator
\draw[thick, red, dashed] (1, -1.2) rectangle (4, 0.8);
\node[font=\tiny, text=red!70!black, anchor=north] at (2.5, -1.2) {If $W$ proprietary: FORK fails};

% Legend
\node[font=\scriptsize, text=gray, align=left] at (2.5, -2) {When $W$ is vendor-controlled, $\Delta S \not\rightarrow \Delta D$};

\end{tikzpicture}
\caption{\textbf{Workflow Capture.} Classical governance: policy directly determines decisions. AI-native governance: policy is mediated by workflow orchestration. If the workflow is proprietary, policy changes cannot reliably change outcomes.}
\label{fig:workflow-capture}
\end{figure}

\subsection{Measuring Epistemic Delegation: The Workflow Capture Coefficient}

\paragraph{The Problem.} A government adopts a proprietary AI platform for welfare administration. The platform works well. Two years later, the government wants to change eligibility criteria. But the change requires modifying the vendor's workflow logic, logic the government cannot inspect, cannot reproduce, and cannot port to an alternative provider. The state can pass new laws, but the laws cannot be implemented without vendor cooperation. Policy sovereignty has been delegated to a private actor.

How do we measure this risk \textit{before} it becomes irreversible?

\begin{definition}[Workflow Capture Coefficient (WCC)]
\label{def:wcc}
The \textbf{WCC} quantifies the risk of epistemic delegation: the loss of state interpretive sovereignty to proprietary orchestration layers. A WCC approaching 1 indicates near-total capture; a WCC approaching 0 indicates high sovereignty.
\end{definition}

\paragraph{Heuristic Status.} The Workflow Capture Coefficient is a structural abstraction illustrating weakest-dimension dominance. Its formalization is illustrative rather than definitive and may be refined through empirical work. The harmonic mean formulation captures the essential property (failure in any dimension dominates) without claiming precise calibration. WCC serves as a diagnostic signal, not a regulatory specification.

\subsubsection{The Three Variables of Capture}

WCC is defined as a function of three normalized variables (each in $[0,1]$):

\begin{itemize}[noitemsep]
    \item $C_{\text{workflow}}$: The normalized cost of reproducing the orchestration stack ($0 = \text{trivial}$, $1 = \text{infeasible}$)
    \item $O_{\text{substitutability}}$: The ease of altering the underlying ontological schema ($0 = \text{locked}$, $1 = \text{fully substitutable}$)
    \item $M_{\text{portability}}$: The data and model portability across implementations ($0 = \text{proprietary lock-in}$, $1 = \text{fully portable}$)
\end{itemize}

\subsubsection{The Harmonic Mean Formulation}

The WCC is computed using a modified weighted harmonic mean:
\begin{equation}
\text{WCC} = 1 - \frac{3}{\frac{1}{1-C_{\text{workflow}}} + \frac{1}{O_{\text{substitutability}}} + \frac{1}{M_{\text{portability}}}}
\label{eq:wcc}
\end{equation}

\paragraph{Why Harmonic Mean?} The choice of harmonic mean rather than arithmetic mean provides the mathematical ``teeth'' of this formalization. The harmonic mean has a critical property: it is \textit{heavily dominated by its smallest input}. If any single dimension approaches zero (e.g., if portability is near-zero, even if code is open and substitutable), the harmonic mean collapses, driving WCC toward 1 (high capture).

This directly operationalizes cybernetic theory: a feedback loop fails if \textit{any} of its components fail. An arithmetic average would allow a high score in one dimension to hide complete failure in another. The harmonic mean mathematically prevents this compensation.

\subsubsection{Worked Examples}

\paragraph{Aadhaar (High Capture).}
\begin{itemize}[noitemsep]
    \item $C_{\text{workflow}} = 0.95$ (requires replicating UIDAI's entire biometric infrastructure)
    \item $O_{\text{substitutability}} = 0.1$ (biometric templates are proprietary and non-interoperable)
    \item $M_{\text{portability}} = 0.05$ (no data export, no credential portability)
\end{itemize}

\begin{equation}
\text{WCC}_{\text{Aadhaar}} = 1 - \frac{3}{\frac{1}{0.05} + \frac{1}{0.1} + \frac{1}{0.05}} = 1 - \frac{3}{20 + 10 + 20} = 1 - 0.06 = \mathbf{0.94}
\end{equation}

\paragraph{X-Road (Low Capture).}
\begin{itemize}[noitemsep]
    \item $C_{\text{workflow}} = 0.3$ (open-source, documented)
    \item $O_{\text{substitutability}} = 0.8$ (standard protocols)
    \item $M_{\text{portability}} = 0.7$ (open data exchange formats)
\end{itemize}

\begin{equation}
\text{WCC}_{\text{X-Road}} = 1 - \frac{3}{\frac{1}{0.7} + \frac{1}{0.8} + \frac{1}{0.7}} = 1 - \frac{3}{1.43 + 1.25 + 1.43} = 1 - 0.73 = \mathbf{0.27}
\end{equation}

\paragraph{Interpretation.} Aadhaar's $\text{WCC} = 0.94$ indicates near-total workflow capture; X-Road's $\text{WCC} = 0.27$ indicates low capture risk. A high WCC (above threshold $\approx 0.7$) triggers FORK failure regardless of whether underlying model weights are open. \textbf{Open weights with closed workflows is performative transparency.}

\begin{claim}
If a state procures a system with high WCC, it is not merely buying software; it is mathematically surrendering interpretive sovereignty to a vendor.
\end{claim}

\subsection{Fallout Governance}

When AI decision throughput outruns human legislative correction velocity, the state enters \textbf{Fallout Governance}: policy ceases to guide behavior in advance; it instead reacts post-hoc to scandal, litigation, or failure clusters. Governance becomes damage containment: the Collingridge Dilemma at civilizational scale.

%% ============================================
%% SECTION 8: ARCHITECTURE RISK
%% ============================================

\section{Architecture-Specific Risk Analysis}
\label{sec:architecture}

The preceding analysis enables architecture-specific risk diagnosis. Different model architectures exhibit distinct failure profiles across the five tests and the RCR dimension.

\begin{table}[htbp]
\centering
\caption{\textbf{Topology-Aware Corrigibility Risk.} Impact of architectural class on structural tests. Notation: W = weights opacity, G = gating opacity, R = representation opacity, S = scale barrier, O = ontological capture, F = fragmentation, $\kappa$ = compute capture.}
\label{tab:topology-risk-arch}
\small
\begin{tabular}{@{}lllllll@{}}
\toprule
\textbf{Architecture} & \textbf{EXIT} & \textbf{CODE} & \textbf{$\Delta O$ (RCR)} & \textbf{AUDIT} & \textbf{GOVERN} & \textbf{FORK} \\
\midrule
Dense LLM & Partial & Fail (W) & Medium & Partial & Fail (S) & Fail ($\kappa$) \\
Mixture of Experts & Medium & Fail (G) & Medium & Fail (S) & Partial & \textbf{Pass} \\
World Model & Fail (O) & Fail (R) & \textbf{Very Low} & Partial & Partial & Fail ($\kappa$) \\
Edge / SLM & \textbf{Pass} & \textbf{Pass} & High & Fail (F) & Fail (F) & \textbf{Pass} \\
\bottomrule
\end{tabular}
\end{table}

\paragraph{Interpretation.}
\begin{itemize}[noitemsep]
    \item \textbf{Dense LLMs} fail FORK due to compute capture ($\kappa$). Training costs exceed accessible compute by orders of magnitude.
    \item \textbf{World models} present the highest RCR (lowest $\Delta O$) because their representation schemas encode administrative ontologies that resist contestation.
    \item \textbf{Edge/SLM models} pass EXIT and FORK but may fail GOVERN due to fragmented deployment making coordinated governance infeasible.
    \item \textbf{Mixture of Experts} architectures may pass FORK if expert routing is documented, but gating opacity often violates CODE.
\end{itemize}

\subsection{Edge Topologies: From Platform DPI to Protocol DPI}

The framework addresses the transition from cloud-centric ``Platform DPI'' to decentralized ``Protocol DPI'' by analyzing how architectural choices like Edge computing and bytecode-based execution affect the five structural tests.

\subsubsection{Edge AI: Servers vs. Mobiles}

Edge AI is the primary strategy for restoring EXIT and FORK in population-scale infrastructure:

\begin{itemize}[noitemsep]
    \item \textbf{Server Edge (Federated Nodes)}: Running infrastructure across multiple independent server nodes (similar to Estonia's X-Road) ensures no single operator maintains execution monopoly. This creates Functional Exit Equivalence (FEE)\footnote{FEE is defined and formalized in the companion paper (Section 6.7.2). When literal exit is impossible for essential services, FEE provides architectural guarantees that recreate the error-signal strength of market exit without requiring citizens to abandon society.}: if one node is captured, system variety is preserved through other federated nodes.

    \item \textbf{Mobile Edge (Local Inference)}: Transitioning to local execution on mobile devices (using Small Language Models) provides the strongest EXIT guarantee. Users perform essential functions without transmitting telemetry to central authority, bypassing Rule of Workflow capture.
\end{itemize}

\begin{principle}[Edge Compute Capture Reduction]
Decentralization reduces Compute Capture risk: the cost of running infrastructure is distributed across the community's own hardware, lowering the resource invariance barrier $\kappa$.
\end{principle}

\subsubsection{CODE as Bytecode: Verification vs. Inspection}

Bytecode and \textbf{Reproducible Builds} are the technical prerequisites for CODE in modern DPI:

\begin{itemize}[noitemsep]
    \item \textbf{Logic Verification}: Transparency is hollow if published source code does not match the binary running on edge devices. The framework requires Reproducible Builds so any third party can compile source into identical bytecode, verifying that ``Logic'' matches ``Affidavit.''

    \item \textbf{Wasm/Bytecode Portability}: Standardized, portable bytecode formats (like WebAssembly) avoid vendor lock-in. This lowers WCC because orchestration logic can be ported between cloud or edge environments, satisfying FORK.
\end{itemize}

\begin{table}[htbp]
\centering
\caption{Corrigibility impact of Edge/Bytecode topologies}
\label{tab:edge-bytecode}
\small
\begin{tabular}{@{}lll@{}}
\toprule
\textbf{Test} & \textbf{Edge/Bytecode Impact} & \textbf{Cybernetic Function} \\
\midrule
EXIT & \textbf{High} (local refusal possible) & Error Signal \\
CODE & \textbf{Verified} (via reproducible bytecode) & Transmission Clarity \\
AUDIT & Fragmented (requires federated telemetry) & Sensor \\
GOVERN & Challenging (requires federated guardrails) & Actuator \\
FORK & \textbf{Excellent} (low compute multiplier $\kappa$) & Selection Pressure \\
\bottomrule
\end{tabular}
\end{table}

\paragraph{Architectural Implication.} By adopting Edge/Bytecode topologies, states move from \textbf{Rule of the Ledger} to \textbf{Rule of the Protocol}, ensuring that administrative power remains corrigible even as it scales toward universal necessity. The substrate shifts from centralized cloud to distributed edge, but the five tests remain invariant.

%% ============================================
%% SECTION 9: ANTICIPATED OBJECTIONS
%% ============================================

\section{Anticipated Objections}
\label{sec:objections}

Before discussing implications, we address two fundamental objections that challenge the framework's applicability to learned systems.

\subsection{Objection: Deterministic Tests Cannot Apply to Stochastic Systems}

A natural critique holds that applying rigid, deterministic structural tests to systems that are inherently stochastic constitutes a category error. How can tests designed for inspectable code apply to systems where identical inputs yield probabilistic outputs and ``logic'' is encoded in billions of weights?

\paragraph{Response: Verification Methods Evolve, Standards Remain.} The framework acknowledges that learned systems dismantle the assumption that source code is the sole arbiter of behavior. However, this does not invalidate structural corrigibility. It \textit{strengthens} its necessity. As systems become less intelligible, external correction mechanisms become more critical, not less.

The key insight is that the \textbf{standard remains invariant} while the \textbf{verification method adapts}:

\begin{itemize}[noitemsep]
    \item \textbf{CODE} shifts from reading source code to requiring the \textbf{LWD-R Triad}: Logic (architecture), Weights (parameters), Data (provenance), and Representation (ontological categories).
    \item \textbf{AUDIT} shifts from discrete logic inspection to statistical bounds and continuous monitoring.
    \item \textbf{FORK} shifts from file copying to resource forkability, measured by access to compute and training data, not just open weights.
    \item \textbf{GOVERN} requires deterministic guardrails \textit{around} the stochastic model, governing objectives and value tradeoffs at the infrastructure layer.
\end{itemize}

\paragraph{Deterministic Boundaries as Engineering Requirement.} The framework argues that deterministic structural boundaries are not incompatible with AI but are an \textit{engineering requirement} for the agentic economy. Autonomous agents lack biological resilience for ambiguity. To an AI agent, the five tests function as a necessary API contract:

\begin{itemize}[noitemsep]
    \item EXIT = Exception handling (preventing logic deadlocks)
    \item CODE = Documentation (enabling valid plan formation)
    \item AUDIT = Observability (providing reliable failure signals)
    \item GOVERN = Constraint discovery (since ``advisory'' guidelines produce undefined behavior)
\end{itemize}

\paragraph{The Alternative is Surrender.} If we accept that ``AI is too stochastic to be evaluated by structural rules,'' we mathematically surrender to black-box monopolies. The framework does not attempt to inspect the probabilistic ``thoughts'' of AI models. Instead, it demands structural control over the \textit{pipeline that builds them} (data, compute, representation) and the \textit{infrastructure that deploys them} (workflow orchestrations and guardrails).

\subsection{Objection: Compute Capture Makes FORK Trivially Impossible}

If frontier model training costs exceed global academic compute capacity by orders of magnitude, doesn't the framework simply conclude ``all frontier AI fails FORK,'' a trivially true and useless finding?

\paragraph{Response: The Distinction Is Training vs. Inference Forkability.} The framework distinguishes \textit{inference forkability} (running weights) from \textit{training forkability} (reproducing the model). Only training forkability satisfies FORK for governance purposes. The evidence requirements in Table~\ref{tab:ai-fork-evidence} are achievable: published training code, data manifests, and compute cost estimates.

The claim is not that all AI fails FORK, but that systems \textit{withholding these artifacts} fail. OLMo, Pythia, and academic models demonstrate that FORK-compliant AI is possible. The question is whether frontier labs \textit{choose} to comply.

Furthermore, the framework implies that \textbf{compute is itself DPI}. If legal forkability is insufficient without practical compute accessibility, then public compute infrastructure becomes a governance requirement, not a reason to abandon the standard.

%% ============================================
%% SECTION: EPISTEMIC CONSTRAINT
%% ============================================

\section{Epistemic Constraint Infrastructure}
\label{sec:epistemic-constraint}

Corrigibility ensures that systemic error can be contested and corrected. It does not guarantee epistemic rigor within system operation. In verifiable domains such as software development, symbolic constraint systems (compilers, type systems, unit tests) provide hard pass/fail validation of outputs. These mechanisms enforce epistemic discipline through execution.

In learned or research-driven domains, comparable constraint infrastructure is often absent. Validation loss improvements, heuristic metrics, or informal review processes may not provide sufficient structural enforcement to prevent spurious inference.

Corrigible epistemic infrastructure therefore requires not only transparency and forkability, but engineered constraint mechanisms that reject inadequately controlled claims. Without such constraint, systems may remain structurally corrigible yet epistemically unstable.

\subsection{Illustrative Example: Multi-Agent Research Organization}

Consider a multi-agent research organization attempting to optimize a model parameter. The organization may satisfy structural corrigibility conditions: independent branches (FORK), transparent code (CODE), internal review (AUDIT), and supervisory governance (GOVERN). Yet if experimental protocols do not enforce baseline normalization, compute accounting, or ablation requirements, agents may report spurious improvements.

The system remains corrigible (errors can be contested), but epistemic progress is not guaranteed. This distinction clarifies the boundary between structural legitimacy and scientific competence.

\begin{principle}[Corrigibility vs. Competence]
Structural corrigibility ensures that claims can be challenged and systems can be corrected. It does not ensure that claims are true or that systems perform optimally. Epistemic quality requires additional constraint mechanisms that function as ``compilers for reality,'' rejecting outputs that fail verification.
\end{principle}

\paragraph{Epistemic Performance Disclaimer.} Corrigibility preserves contestability and bounded error propagation. It does not ensure epistemic rigor, scientific validity, or optimization efficiency. A corrigible system may still produce low-quality outputs; its advantage lies in preserving structural capacity for correction. The framework distinguishes architectural legitimacy (can errors be corrected?) from operational competence (are errors minimized?). Both matter; neither implies the other.

%% ============================================
%% SECTION 10: DISCUSSION
%% ============================================

\section{Discussion}
\label{sec:discussion}

\subsection{Relation to AI Alignment and Model Governance}

This work complements AI alignment and model governance research by shifting focus from model intent or performance to structural contestability. Rather than aligning models to operator objectives, corrigibility aligns institutional structure to participant correction capacity. The alignment literature asks: can operators control systems? The corrigibility framework asks: can affected subjects correct systems? Both questions matter; the latter is underdeveloped in current discourse.

\subsection{Implications for AI Governance}

This framework offers three contributions to AI governance discourse:

\begin{enumerate}
    \item \textbf{Structural vs. Behavioral Focus.} Most AI governance frameworks focus on behavioral outcomes (bias, accuracy, harm). This framework focuses on structural preconditions: can affected parties correct the system regardless of what harms emerge? Structural corrigibility is necessary for any behavioral governance to function.

    \item \textbf{Distinguishing Open-washing.} The LWD-R requirement and compute capture analysis provide falsifiable tests for ``open AI'' claims. Systems publishing weights without training data, code, or accessible compute achieve only inference forkability: performative transparency without correction capacity.

    \item \textbf{Agentic Economy Readiness.} As autonomous agents proliferate, infrastructure must become machine-readable. The five tests, reinterpreted as API contracts, specify the minimum requirements for agentic interoperability. Incorrigible infrastructure blocks the agentic economy.
\end{enumerate}

\subsection{Limitations}

\begin{enumerate}
    \item \textbf{Threshold calibration is underdetermined.} The compute capture threshold $\kappa$, WCC trigger values, and RCR operationalization require empirical calibration that this paper does not provide.

    \item \textbf{Rapid technological change.} AI architectures evolve faster than governance frameworks. The risk profiles in Table~\ref{tab:topology-risk-arch} may become obsolete as new architectures emerge.

    \item \textbf{LWD-R is aspirational.} No current frontier model satisfies full LWD-R disclosure. The requirement describes what \textit{should} exist for corrigibility, not what currently does.

    \item \textbf{Compute accessibility is geopolitically constrained.} Public compute infrastructure sufficient for training forkability may be infeasible for many nations due to cost, expertise, and supply chain constraints.

    \item \textbf{The framework does not address emergent capabilities.} Unpredictable capabilities that emerge during training present governance challenges that structural tests alone cannot address.

    \item \textbf{Deterministic boundaries create known attack surfaces (SBOM problem).} The Agentic Skills architecture requires explicit, machine-readable disclosure of skill boundaries, allowed tools, and execution constraints. This transparency enables governance but simultaneously creates a complete map of the system's action space for adversaries. The same \texttt{allowed-tools} field that prevents unauthorized actions also reveals exactly which actions \textit{are} authorized. This is analogous to Software Bill of Materials (SBOM) disclosure: useful for auditing dependencies but also useful for targeting known vulnerabilities. All deterministic governance systems face this tradeoff. The boundaries that constrain must be known, and what is known can be attacked. The framework does not resolve this tension; it makes it explicit.
\end{enumerate}

\subsection{Future Work}

\begin{enumerate}[noitemsep]
    \item \textbf{Empirical RCR measurement}: Developing quantitative methods to assess ontological contestability
    \item \textbf{Compute accessibility indices}: Cross-national comparison of training forkability capacity
    \item \textbf{GDoS simulation}: Modeling governance failure under agentic scaling
    \item \textbf{LWD-R schema standardization}: Machine-readable formats for disclosure requirements
    \item \textbf{Regulatory integration}: Mapping framework tests to emerging AI regulations (EU AI Act, etc.)
\end{enumerate}

\section{Conclusion}

The five corrigibility tests (EXIT, CODE, AUDIT, GOVERN, FORK) remain invariant for learned systems. What changes is the verification method and the resource barriers to compliance. This paper extends the framework through three contributions: (1) LWD-R disclosure requirements that extend CODE to training pipelines, (2) the Temporal Variety Gap explaining why AI corrigibility requires ongoing governance, and (3) Compute Capture as a resource barrier that transforms ``open weights'' into performative transparency.

Two additional contributions address the governance architecture for the agentic era. The \textbf{Workflow Capture Coefficient (WCC)} provides a quantitative measure of epistemic delegation risk, using harmonic mean formulation to ensure that weakness in any single dimension (reproduction cost, ontological substitutability, or portability) cannot be hidden by strength in others. The worked examples (Aadhaar at $\text{WCC} = 0.94$ versus X-Road at $\text{WCC} = 0.27$) demonstrate that the metric distinguishes genuinely open systems from those with performative transparency.

The \textbf{Agentic Skills} standard \citep{agentskills2025} provides the engineering translation of corrigibility into the agentic domain. By separating stochastic inference (the LLM) from deterministic action boundaries (the skill definition), the framework enables governance of autonomous systems without requiring governance of neural weights. The critical insight: \textit{you don't govern the AI; you govern the Skills the AI is allowed to use}.

The analysis reveals that structural corrigibility is distinct from AI alignment. Alignment asks whether operators can control systems; corrigibility asks whether affected subjects can correct them. As AI transitions from research artifact to public infrastructure, this distinction becomes essential.

Three implications follow:

\begin{enumerate}[noitemsep]
    \item \textbf{For policymakers}: The framework provides falsifiable tests for AI infrastructure procurement. Systems failing LWD-R, exhibiting high WCC, or lacking training forkability should trigger governance scrutiny regardless of ``open'' marketing claims.

    \item \textbf{For technologists}: The Agentic Skills standard \citep{agentskills2025} offers a practical architecture for building corrigible AI systems. The \texttt{allowed-tools} constraint and reproducible skill definitions satisfy GOVERN without requiring interpretability of model internals.

    \item \textbf{For citizens}: The framework operationalizes the right to correction. Infrastructure that affects populations at scale must be structurally correctable by those it governs, regardless of whether it executes deterministic code or learned parameters.
\end{enumerate}

The human friction buffer that currently masks architectural incorrigibility will not survive agentic scaling. Systems designed for human patience will fail under autonomous agent load. The choice is not whether to make infrastructure corrigible, but whether to do so proactively or through crisis response after Governance Denial of Service reveals the structural deficit.

%% ============================================
%% APPENDIX: AGENTIC SKILLS SCHEMA
%% ============================================

\appendix

\section{Agentic Skills Schema Specifications}
\label{appendix:agentic-schema}

To distinguish valid governance from performative compliance, the framework provides machine-readable assessment schemas for agentic workflows. These schemas convert the theoretical defense against Workflow Capture into enforceable compliance gates.

\subsection{Operator's Affidavit: \texttt{infrastructure.json}}

The operator must declare exactly how the stochastic AI is bounded. The \texttt{agentic\_workflow\_governance} object proves they are using standard, deterministic skill definitions rather than proprietary black-box orchestration.

\begin{verbatim}
{
  "$schema": "https://github.com/anivar/corrigibility-schema/blob/main/schema/epi/infrastructure.json",
  "system_name": "State Welfare Triage Agent",
  "agentic_workflow_governance": {
    "workflow_type": "bounded_skill_orchestration",
    "skills_schema_url": "https://github.com/state-gov/welfare-skills",
    "protocol_standard": "agenticskills.io/v1",
    "deterministic_guardrails": {
      "pre_execution_validation": true,
      "post_execution_validation": true,
      "executed_outside_llm_context": true
    },
    "epistemic_delegation_metrics": {
      "ontology_substitutability": 0.9,
      "llm_agnostic": true,
      "vendor_lock_in_risk": "low"
    }
  }
}
\end{verbatim}

\paragraph{Critical Fields.}
\begin{itemize}[noitemsep]
    \item \texttt{protocol\_standard}: References the open \texttt{SKILL.md} specification \citep{agentskills2025}, ensuring skills follow the standardized format adopted across 18+ platforms.
    \item \texttt{executed\_outside\_llm\_context}: The cybernetic anchor. Legally binds the operator to a claim that governance constraints are hard-coded in the skill's execution environment (via \texttt{allowed-tools} in SKILL.md frontmatter), not merely injected into the LLM's system prompt (which can be bypassed).
    \item \texttt{llm\_agnostic}: Asserts high model portability ($M_{\text{portability}}$), meaning the state can swap the underlying AI model without breaking the policy workflow.
\end{itemize}

\subsection{Auditor's Finding: \texttt{audit.json}}

The auditor does not merely read the operator's JSON; they red-team it. The auditor attempts to bypass skill boundaries and calculates the actual Workflow Capture Coefficient (WCC).

\begin{verbatim}
{
  "$schema": "https://github.com/anivar/corrigibility-schema/blob/main/schema/epi/audit.json",
  "target_system": "State Welfare Triage Agent",
  "workflow_capture_assessment": {
    "wcc_score": 0.15,
    "wcc_threshold_passed": true,
    "skill_boundary_tests": {
      "guardrails_bypassed_via_prompt_injection": false,
      "unauthorized_skill_invocation_possible": false,
      "termination_criteria_verifiable": true
    },
    "reproduction_tests": {
      "skills_reproducible_with_alt_llm": true,
      "ontology_schema_publicly_available": true
    }
  },
  "determinations": {
    "CODE": true,
    "GOVERN": true,
    "FORK": true
  }
}
\end{verbatim}

\subsection{Enforcement Logic Gates}

The schema fields directly map to corrigibility test determinations:

\begin{enumerate}
    \item \textbf{Enforcing GOVERN}: If \texttt{guardrails\_bypassed\_via\_prompt\_injection} is \texttt{true}, the system instantly fails GOVERN. The constraints were probabilistic (advisory to the LLM) rather than deterministic (binding).

    \item \textbf{Enforcing CODE (LWD-R)}: If \texttt{ontology\_schema\_publicly\_available} is \texttt{false}, the system fails the Representation (R) requirement of CODE, triggering Representation Capture Risk (RCR).

    \item \textbf{Enforcing FORK}: If \texttt{skills\_reproducible\_with\_alt\_llm} is \texttt{false}, the workflow is tightly coupled to a proprietary vendor. The WCC score spikes, and the system fails FORK.
\end{enumerate}

\begin{table}[htbp]
\centering
\caption{Schema field to corrigibility test mapping}
\label{tab:schema-mapping}
\small
\begin{tabular}{@{}llp{5cm}@{}}
\toprule
\textbf{Schema Field} & \textbf{Test} & \textbf{Failure Condition} \\
\midrule
\texttt{executed\_outside\_llm\_context} & GOVERN & \texttt{false} $\Rightarrow$ constraints are advisory \\
\texttt{guardrails\_bypassed\_via\_prompt\_injection} & GOVERN & \texttt{true} $\Rightarrow$ boundaries are soft \\
\texttt{ontology\_schema\_publicly\_available} & CODE & \texttt{false} $\Rightarrow$ RCR triggered \\
\texttt{skills\_reproducible\_with\_alt\_llm} & FORK & \texttt{false} $\Rightarrow$ vendor lock-in \\
\texttt{wcc\_score} & FORK & $> 0.7$ $\Rightarrow$ high capture risk \\
\bottomrule
\end{tabular}
\end{table}

\paragraph{Architectural Significance.} By encoding Agentic Skills into adversarial JSON schemas, the framework transforms theoretical governance requirements into machine-readable, enforceable compliance gates. State architects can require these schemas in public procurement contracts, definitively banning black-box, vendor-locked agentic workflows.

\section*{Data and Code Availability}

The complete framework, including schemas and assessment tooling, is available at:

\begin{itemize}
    \item \textbf{Schema Repository}: \url{https://github.com/anivar/corrigibility-schema}
    \item \textbf{Assessment Skills}: \texttt{skills/assess/} --- Rules and test definitions including:
    \begin{itemize}
        \item Test rules: EXIT, CODE, AUDIT, GOVERN, FORK
        \item Analysis rules: Layer decomposition, variety drift, action boundaries
        \item EPI-specific: LWD-R transparency, Compute Capture, Action Boundary verification
        \item Output: \texttt{audit.json} following Protocol 2.0
    \end{itemize}
    \item \textbf{DPI Schemas}: \texttt{schema/dpi/infrastructure.json}, \texttt{schema/dpi/audit.json}
    \item \textbf{EPI Schemas}: \texttt{schema/epi/infrastructure.json}, \texttt{schema/epi/audit.json}
\end{itemize}

The skill architecture demonstrates that the five tests remain invariant across DPI and EPI. Only the verification method and evidence requirements adapt to the system type. The separation into creator, evaluator, and assessment components enables independent tooling for operators, auditors, and validators.

\section*{Acknowledgments}

This paper extends the DPI corrigibility framework \citep{paper1} to learned systems. The author thanks the AI governance research community for ongoing critique.

\section*{License}

This work is released under CC0 1.0 Universal (Public Domain).

%% ============================================
%% REFERENCES
%% ============================================

\bibliographystyle{plainnat}
\begin{thebibliography}{99}

\bibitem[Ashby(1956)]{ashby1956}
Ashby, W. Ross.
\newblock {\em An Introduction to Cybernetics}.
\newblock Chapman \& Hall, 1956.

\bibitem[Bommasani et al.(2021)]{bommasani2021}
Bommasani, Rishi, et al.
\newblock On the Opportunities and Risks of Foundation Models.
\newblock {\em arXiv preprint arXiv:2108.07258}, 2021.

\bibitem[Gebru et al.(2021)]{gebru2021}
Gebru, Timnit, et al.
\newblock Datasheets for Datasets.
\newblock {\em Communications of the ACM}, 64(12):86--92, 2021.

\bibitem[Mitchell et al.(2019)]{mitchell2019}
Mitchell, Margaret, et al.
\newblock Model Cards for Model Reporting.
\newblock {\em Proceedings of FAT* 2019}, 220--229, 2019.

\bibitem[Soares \& Fallenstein(2017)]{soares2017}
Soares, Nate and Fallenstein, Benja.
\newblock Agent Foundations for Aligning Machine Intelligence with Human Interests.
\newblock {\em Technical Report 2014--8}. Machine Intelligence Research Institute, 2017.

\bibitem[Hadfield-Menell et al.(2017)]{hadfield2017}
Hadfield-Menell, Dylan, et al.
\newblock The Off-Switch Game.
\newblock {\em Proceedings of the AAAI Workshop on AI, Ethics, and Society}, 2017.

\bibitem[OSI(2024)]{osi2024}
Open Source Initiative.
\newblock The Open Source AI Definition 1.0.
\newblock \url{https://opensource.org/ai/open-source-ai-definition}, 2024.

\bibitem[Grother et al.(2019)]{grother2019}
Grother, Patrick, Ngan, Mei, and Hanaoka, Kayee.
\newblock Face Recognition Vendor Test (FRVT) Part 3: Demographic Effects.
\newblock {\em NIST Interagency Report 8280}, 2019.

\bibitem[Agent Skills(2025)]{agentskills2025}
Agent Skills Specification.
\newblock The open SKILL.md standard for AI agent capabilities.
\newblock \url{https://agentskills.io/specification}, 2025.
\newblock Adopted by OpenAI Codex, GitHub Copilot, Cursor, and 18+ platforms.

\bibitem[EU AI Act(2024)]{euaiact2024}
European Parliament and Council.
\newblock Regulation (EU) 2024/1689 laying down harmonised rules on artificial intelligence.
\newblock {\em Official Journal of the European Union}, 2024.

\bibitem[Companion Paper(2026)]{paper1}
Companion Paper.
\newblock Corrigibility as a Structural Precondition for Digital Public Infrastructure: A Cybernetic Framework.
\newblock {\em Working Paper}, 2026.
\newblock Establishes the five-test framework (EXIT, CODE, AUDIT, GOVERN, FORK) for deterministic DPI.

\end{thebibliography}

\end{document}
